% Generated draft: requires mathematical and source review.
\documentclass[11pt]{article}
\usepackage[margin=25mm]{geometry}
\usepackage{fontspec}
\setmainfont{Latin Modern Roman}
\usepackage{amsmath,amssymb,mathtools,unicode-math}
\setmathfont{Latin Modern Math}
\usepackage{enumitem,needspace,xurl,hyperref}
\setlength{\parindent}{0pt}
\setlength{\parskip}{5pt}
\setlength{\emergencystretch}{3em}
\newcommand{\N}{\mathbb{N}}
\newcommand{\Z}{\mathbb{Z}}
\newcommand{\Q}{\mathbb{Q}}
\newcommand{\R}{\mathbb{R}}
\newcommand{\C}{\mathbb{C}}
\newcommand{\F}{\mathbb{F}}
\newcommand{\A}{\mathbb{A}}
\newcommand{\PP}{\mathbb{P}}
\newcommand{\E}{\mathbb{E}}
\newcommand{\eps}{\varepsilon}
\newcommand{\dd}{\,\mathrm{d}}
\begin{document}
\section*{UnsolvedMath ID: 800142}
\textbf{Source problem code:} catalogue entry
\subsection{Definitions and mathematical statement}
Let \(q\) be a positive integer, and let \(a,b\in\Z/q\Z\) be distinct reduced residue classes, meaning that each class is represented by an integer coprime to \(q\). For real \(x\geq 2\), define
\[
\pi(x;q,a)=\#\{p\leq x:p\text{ is prime and }p\equiv a\pmod q\}.
\]
Define the set of ties by
\[
T_{q;a,b}=\{x\in[2,\infty):\pi(x;q,a)=\pi(x;q,b)\}.
\]
For a measurable set \(S\subseteq[2,\infty)\), its natural and logarithmic densities, when the respective limits exist, are
\[
d(S)=\lim_{X\to\infty}\frac{1}{X}\int_2^X\mathbf{1}_S(t)\,\dd t,
\qquad
\delta(S)=\lim_{X\to\infty}\frac{1}{\log X}\int_2^X\mathbf{1}_S(t)\,\frac{\dd t}{t}.
\]
\textbf{Problem.} For every such \(q,a,b\), prove unconditionally that \(T_{q;a,b}\) has density zero. The source explicitly allows natural or logarithmic density: the requested conclusion is \(d(T_{q;a,b})=0\), or alternatively \(\delta(T_{q;a,b})=0\), including existence of the relevant limit. The displayed definitions use real values of \(x\), as is conventional for the prime-counting function; the source does not explicitly specify whether its density is taken over real or integer values.

\textbf{Qualification concerning the supplied conditional background.} The source invokes the generalized Riemann hypothesis (GRH) and a very mild version of a linear independence hypothesis (LI), referring to [Devin, Theorem 2.2]. It asserts existence and absolute continuity with respect to Lebesgue measure of a limiting logarithmic distribution, followed by a stronger conclusion for every function \(f(x)=o(\sqrt{x}/\log x)\). However, the two expressions printed in that discussion are
\[
\bigl(\pi(x;q,a)=\pi(x;q,b)\bigr)\frac{\log x}{\sqrt{x}}
\quad\text{and}\quad
\bigl|\pi(x;q,a)=\pi(x;q,b)\bigr|<f(x).
\]
These contain equality relations where numerical quantities are required. Replacing those equality signs by subtraction would produce a normalized prime-counting discrepancy and a small-discrepancy condition, but that repair is an interpretation requiring source confirmation. The exact LI hypothesis and the full cited theorem are also absent. Accordingly, this conditional background is retained as a qualified source assertion, not as a precisely established auxiliary theorem.
\subsection{Short English statement}
For every modulus and every pair of distinct reduced residue classes, prove without unproved hypotheses that the values of \(x\) at which the two prime counts tie have natural or logarithmic density zero. The source's stronger conditional discussion contains malformed expressions and an unspecified LI hypothesis, which require review.
\subsection{Sources}
\begin{itemize}
\item Local UnsolvedMath registry, permanent ID \(800142\), exact source code \texttt{catalogue entry}, titled ``Comparative Prime Number Theory --- Problem 11.'' The supplied registry labels the problem open; this label is not an independently verified assessment.
\item Hamieh, Kadiri, Martin, and Ng, \emph{Comparative Prime Number Theory Problem List} (2024), Problem (\#11), \url{https://arxiv.org/abs/2407.03530}. The supplied provenance records access on 2026-10-16 and extraction by \texttt{arxiv\_latex}; no independent source visit or verification is asserted here.
\item The source identifies Greg Martin, University of British Columbia, as the contributor, attributes the question to Carl Pomerance, and says that it appears in [Mar, Section 3]. No full bibliographic entry for [Mar] is supplied.
\item The conditional background cites [Devin, Theorem 2.2] for a mild LI assumption. No full bibliographic entry, link, or theorem statement for this reference is supplied.
\end{itemize}
\end{document}
